\chapter{增强版的协程调度器} \label{ch:coroutine-scheduler}

第~\ref{ch:simple-task}~章实现的 \texttt{simple\_task} 是一个“无栈”的协作式调度器：定时器中断负责计时，主循环负责执行任务，任务回调每次运行后都会返回，下一次运行再从回调函数的入口重新开始。这种模型简单、可靠，但它有一个天然的局限：任务无法“运行到一半时暂停，等条件满足后再从暂停的地方继续”。

例如，下面的代码在普通函数里是理所当然的：

\begin{lstlisting}[style=c]
void SensorTask(void)
{
    Sensor_Start();        /* 启动一次传感器转换 */
    WaitConvertDone();     /* 希望在这里暂停，转换完成后继续 */
    value = Sensor_Read();
    SendResult(value);
}
\end{lstlisting}

在 \texttt{simple\_task} 里，如果不想使用状态机，就只能用“延时返回”模拟这种效果。但一旦任务的暂停点藏在一个多层的函数调用里，或者暂停时需要保留大量局部变量，无栈模型就会变得非常别扭。

本章的“增强版协程调度器”换了一种思路：让任务在任意位置真正暂停，把 CPU 交给其他任务，之后还能精确回到原来的暂停点继续执行。它的实现基础是 C 标准库的 \texttt{setjmp} 与 \texttt{longjmp}，配合为每个任务准备的独立栈（或共享栈备份），把一个普通 C 函数变成可以“挂起和恢复”的协程。

\medskip 本章涉及到的实验内容：
\begin{enumerate}
\item setjmp 与 longjmp 上下文跳转实验；
\item 独立栈协程任务创建与恢复实验；
\item Task\_Wait 定时挂起实验；
\item 共享栈任务与栈备份实验；
\item Task\_Join 与条件变量实验。
\end{enumerate}

完整代码位于 \texttt{Chapter-coroutine-scheduler/协程调度器setjmp/} 目录，包括主文件 \texttt{task\_scheduler.c} 与 \texttt{task\_scheduler.h}、配置文件\texttt{task\_scheduler\_config.h} 和演示用的 \texttt{main.c}。

% ---------------------------------------------------------------------
\section{setjmp 与 longjmp：C 语言里的“现场保存和恢复”}

\texttt{setjmp} 与 \texttt{longjmp} 是 C 标准库提供的非局部跳转工具，头文件为 \texttt{<setjmp.h>}。它们能让程序保存“当前执行现场”，并在之后的任意时刻恢复这个现场，使控制流回到保存现场的位置。

\begin{itemize}
\item \texttt{setjmp(jmp\_buf env)}：把当前函数调用现场（寄存器、栈指针、程序计数器等）保存到 \texttt{env}，第一次调用时返回 0；
\item \texttt{longjmp(jmp\_buf env, int val)}：恢复 \texttt{env} 保存的现场，并让对应位置的 \texttt{setjmp} 以 \texttt{val} 作为返回值返回，程序就像“从刚才暂停的地方继续”一样。
\end{itemize}

一个最直观的例子：

\begin{lstlisting}[style=c]
#include <stdio.h>
#include <setjmp.h>

jmp_buf env;

int main(void)
{
    if (setjmp(env) == 0) {
        /* 第一次到达这里 */
        printf("第一次运行\n");
        longjmp(env, 1);   /* 跳回 setjmp，并让它返回 1 */
    } else {
        /* longjmp 之后，程序回到这里继续执行 */
        printf("从 longjmp 恢复\n");
    }
    return 0;
}
\end{lstlisting}

从表面看，\texttt{setjmp} 只“执行”了一次，却可能“返回”两次：第一次返回 0，表示现场刚被保存；之后每次 \texttt{longjmp} 回到该位置时，都以非零值返回，表示“现场被恢复了”。

\begin{anybox}{使用 setjmp 与 longjmp 的注意事项}

\begin{enumerate}
\item \texttt{setjmp} 保存的是“CPU 现场”，不是“整个内存”。局部变量若在 \texttt{setjmp} 之后被修改，恢复后这些修改是否保留取决于编译器实现，因此跨 \texttt{setjmp} 使用的局部变量建议声明为 \texttt{volatile}；
\item \texttt{jmp\_buf} 只在保存现场的函数仍处于活动状态时有效。如果函数已经返回，再对它的 \texttt{jmp\_buf} 调用 \texttt{longjmp} 属于未定义行为；
\item \texttt{longjmp} 不是普通函数调用，它不会“展开”栈帧，也不会调用析构和清理逻辑，因此不能代替 C++ 异常或 \texttt{return}；
\item 恢复现场时，栈指针会变回保存时刻的值。这也是为什么协程必须认真管理“每个任务的栈”。
\end{enumerate}

\end{anybox}
% ---------------------------------------------------------------------
\section{为什么栈空间如此重要}

一个正在运行的 C 函数，其状态并不只保存在“变量”里，还保存在调用栈上：

\begin{itemize}
\item 函数每调用一层，就把返回地址、参数、局部变量压栈；
\item 函数返回时，再把这些数据从栈上弹出；
\item 栈指针（SP）指向当前栈顶，它决定了“下一条返回地址在哪里”。
\end{itemize}

因此，“暂停一个任务”并不仅仅意味着保存几个变量的值，而是要把整个调用栈的状态保存下来。只要任务是在 \texttt{TaskFunc1} 的某次循环里暂停，暂停点之上可能还有 \texttt{printf}、传感器库函数等好几层栈帧，每一层都包含局部变量和返回地址。恢复任务时，这些栈帧必须原封不动地回来，函数才能“像什么都没发生一样”继续执行。

\begin{anybox}{栈与暂停深度}
任务的“暂停点”越深，需要保存的调用栈帧就越多，这正是协程调度器必须为每个任务准备独立栈（或共享栈备份）的根本原因。
\end{anybox}

如果所有任务都使用同一个栈，会发生什么？

\begin{enumerate}
\item 任务 A 调用到第三层函数后暂停，此时栈上有 A 的完整调用链；
\item 调度器切换给任务 B，B 开始压栈，A 的栈帧被新数据覆盖；
\item 等 A 再次被调度时，它的栈内容已经被破坏，程序必然崩溃。
\end{enumerate}

解决办法有两种：

\begin{itemize}
\item 每个任务拥有独立栈：任务暂停后，它的栈原封不动地留在自己的内存区域，切换任务只需把栈指针切到另一个任务的栈顶即可；
\item 共享栈 + 备份区：多个任务共用同一个运行栈，任务暂停前把当前栈内容整体拷贝到自己的备份缓冲区，恢复时再拷贝回去。这种方法节省 RAM，但每次切换都要付出拷贝成本，且必须精确计算实际使用的栈大小。
\end{itemize}

本章的调度器两种模式都支持：\texttt{Task\_CreateWithStack} 创建独立栈任务， \texttt{Task\_Create} 创建共享栈任务。

% ---------------------------------------------------------------------
\subsection{独立栈任务}

独立栈任务的示意图如图~\ref{fig:unique-stack}~所示：每个任务拥有一段互不重叠的静态数组作为自己的栈，调度器切换任务时改变栈指针即可。

\begin{lstlisting}[style=c]
/* 任务栈（独立栈）：每个任务分配自己的内存 */
static uint8_t s_Stack1[1024 * 10];
static uint8_t s_Stack2[1024 * 10];

Task_CreateWithStack(TaskFunc1, "task1", s_Stack1, sizeof(s_Stack1), 10);
Task_CreateWithStack(TaskFunc2, "task2", s_Stack2, sizeof(s_Stack2), 20);
\end{lstlisting}

\begin{figure}[htbp]
\centering
    \begin{tikzpicture}[
        stack/.style={draw=structurecolor, thick, minimum width=3.6cm,
                      minimum height=4.8cm, align=center, font=\small},
        lab/.style={font=\small}
    ]
        \node[stack] (s1) at (0,0) {任务 1 栈\\[4pt]
            \texttt{s\_Stack1[10240]}};
        \node[stack] (s2) at (6.2,0) {任务 2 栈\\[4pt]
            \texttt{s\_Stack2[10240]}};
		\node[lab, align=center] at (0,-2.9) {任务1暂停后\\栈内容保留};
		\node[lab, align=center] at (6.2,-2.9) {任务2运行\\使用自己的栈};
        \node[lab] at (-2.7,2.2) {\texttt{stackTop}};
        \draw[-stealth, thick] (-2.2,1.8) -- (-1.0,1.8);
        \node[lab] at (3.5,2.2) {\texttt{stackTop}};
        \draw[-stealth, thick] (4.0,1.8) -- (5.2,1.8);
    \end{tikzpicture}
\caption{独立栈任务：每个任务拥有自己的栈空间} \label{fig:unique-stack}
\end{figure}

独立栈的优点是恢复速度最快，缺点是 RAM 开销大。任务栈大小必须大于该任务“最坏调用深度”下实际占用的栈，否则会发生栈溢出，悄悄破坏相邻内存。

% ---------------------------------------------------------------------
\subsection{共享栈任务与栈备份}

共享栈任务没有自己的运行栈，而是在任务暂停前把“当前栈上用过的部分”备份到私有缓冲区：

\begin{lstlisting}[style=c]
static void Task_BackupSharedStack(TaskControlBlock_T *pTCB)
{
    /* 计算从栈顶到当前栈指针之间实际使用了多少字节 */
    pTCB->stackUsed = s_TaskScheduler.pfnCalcStackUsage(pTCB);

    /* 若备份空间不足，死循环提示增大配置 */
    while (pTCB->stackUsed > TASK_SHARED_BAK_STACK_SIZE);

    /* 把用过的栈内容拷贝到任务自己的备份区 */
    Task_MemCopy(pTCB->pBackupStack,
                 (uint8_t *)(pTCB->stackTop - pTCB->stackUsed),
                 pTCB->stackUsed);
}
\end{lstlisting}

共享栈的优点是多个任务只需少量备份空间，RAM 占用小；缺点是每次挂起和恢复都要执行一次内存拷贝，切换成本更高，而且要求所有共享栈任务的“实际最大栈用量”之和小于一个备份区能容纳的上限。这正是配置项 \texttt{TASK\_SHARED\_BAK\_STACK\_SIZE} 存在的原因。

% ---------------------------------------------------------------------
\section{任务控制块：记录“现场”与“栈”}

调度器为每个任务维护一个任务控制块（TCB），本章的 \texttt{TaskControlBlock\_T} 比 \texttt{simple\_task} 的任务池多了两类关键信息：

\begin{enumerate}
\item 栈信息：\texttt{stackTop}、\texttt{stackSize}、\texttt{stackType}，记录任务栈在哪里、有多大；
\item 跳转现场：\texttt{jmp\_buf env}，记录任务被挂起时的 CPU 现场；
\item 唤醒条件：\texttt{nextRunTime}（定时唤醒）和 \texttt{pWaitCondition}（条件唤醒）。
\end{enumerate}

\begin{lstlisting}[style=c]
typedef struct __TaskControlBlock_T {
    char name[TASK_NAME_LENGTH];
    uint16_t pid;
    uint8_t state;                  /* 任务状态 */
    TaskCondition_T condition;      /* 任务退出条件 */

    size_t stackTop;                /* 栈顶地址（高地址） */
    size_t stackSize;               /* 独立栈总大小 */
    TaskStackType_T stackType;      /* 独立栈 or 共享栈 */

    uint32_t nextRunTime;           /* 下次唤醒的绝对时间（ms） */
    int entryArg;
    jmp_buf env;                    /* 任务挂起现场 */
    TaskEntry_f pfnEntry;
    TaskCondition_T *pWaitCondition;

#if (TASK_MAX_SHARED_TASK > 0)
    uint8_t *pBackupStack;          /* 共享栈备份缓冲区 */
    size_t stackUsed;               /* 实际使用的栈大小 */
#endif
    struct __TaskControlBlock_T *pNext;
} TaskControlBlock_T;
\end{lstlisting}

可以看到，TCB 是“任务控制信息”和“任务执行现场”的桥梁：状态机决定任务什么时候能跑， \texttt{jmp\_buf} 决定跑起来后从哪里继续，栈信息则保证恢复现场时有地方可回。

% ---------------------------------------------------------------------
\section{调度器的状态机}

本章调度器把任务抽象成 5 个状态，如表~\ref{tab:coro-state}~所示。

\begin{table}[htbp]
\centering \caption{增强版协程调度器的任务状态} \label{tab:coro-state}
    \begin{tabular}{llp{8cm}}
        \toprule
        状态 & 含义 & 说明 \\
        \midrule
        \texttt{INIT}     & 初始态 & 任务刚创建，尚未进入入口函数 \\
        \texttt{READY}    & 就绪态 & 时间到或条件满足，等待被选中执行 \\
        \texttt{RUNNING}  & 运行态 & 正在使用 CPU 的任务 \\
        \texttt{SUSPEND}  & 挂起态 & 任务调用了 \texttt{Task\_Wait} 或条件等待，正在休眠 \\
        \texttt{DELETED}  & 结束态 & 任务入口函数返回，等待回收 \\
        \bottomrule
    \end{tabular}
\end{table}

状态转移可以概括为如图\ref{fig:coro-state}所示的那样。

\begin{figure}[htbp]
	\centering
	\begin{tikzpicture}[
		scale=0.7,                % 整体缩放比例
		transform shape,          % 节点文字、箭头等一并缩放
		state/.style={
			rectangle, draw=structurecolor, thick, rounded corners=3pt,
			minimum width=2.4cm, minimum height=0.9cm, align=center,
			fill=structurecolor!8, font=\small
		},
		arrow/.style={-stealth, thick, structurecolor},
		label/.style={font=\footnotesize\itshape, align=center}
		]
		
		% 状态节点
		\node[state] (init)     {INIT};
		\node[state, right=2.8cm of init] (running) {RUNNING};
		\node[state, right=2.8cm of running] (suspend) {SUSPEND};
		\node[state, right=2.8cm of suspend] (ready) {READY};
		\node[state, below=1.0cm of running] (deleted) {DELETED};
		
		% 转换箭头
		\draw[arrow] (init) -- node[label, above] {调度器首次启动} (running);
		\draw[arrow] (running) -- node[label, above] {调用\\\texttt{Task\_Wait} 或\\条件不满足} (suspend);
		\draw[arrow] (suspend) -- node[label, above] {超时到期\\或条件置位} (ready);
		\draw[arrow] (ready) to[out=45, in=135]
		node[label, above, sloped] {调度器选择} (running);
		\draw[arrow] (running) -- node[label, left] {任务函数返回} (deleted);
		
		% 可选：从 INIT 直接到 DELETED（通常不需要，省略）
		
	\end{tikzpicture}
	\caption{协程调度器任务状态转换图}
	\label{fig:coro-state}
\end{figure}

其中：

\begin{itemize}
\item \texttt{INIT} $\rightarrow$ \texttt{RUNNING}：调度器首次启动该任务；
\item \texttt{RUNNING} $\rightarrow$ \texttt{SUSPEND}：任务调用 \texttt{Task\_Wait} 与 \texttt{Task\_ConditionWait}，主动让出 CPU；
\item \texttt{SUSPEND} $\rightarrow$ \texttt{READY}：定时时间到或条件被置位；
\item \texttt{READY} $\rightarrow$ \texttt{RUNNING}：调度器选中该任务并 \texttt{longjmp} 回到它的现场；
\item \texttt{RUNNING} $\rightarrow$ \texttt{DELETED}：任务入口函数正常返回。
\end{itemize}

% ---------------------------------------------------------------------
\section{状态机“被简化”在哪里}

\texttt{simple\_task} 的状态机写在调度器的数据结构和循环里；而本章协程的状态机，很大一部分被“藏”进了 \texttt{setjmp} 与 \texttt{longjmp}。任务代码本身不需要手写 \texttt{switch(state)} 来记住自己执行到哪一步。

以 \texttt{Task\_Wait} 为例，任务“暂停”和“恢复”只需要几行：

\begin{lstlisting}[style=c]
void Task_Wait(uint32_t timeMs)
{
    /* 1. 记录唤醒时间，进入挂起态 */
    s_TaskScheduler.pCurrent->nextRunTime =
        s_TaskScheduler.pfnGetTime() + timeMs;
    s_TaskScheduler.pCurrent->state = TASK_STATUS_SUSPEND;

    /* 2. 保存当前任务现场；若被恢复，setjmp 会返回 TASK_COMMON_RUN */
    if (setjmp(s_TaskScheduler.pCurrent->env) != TASK_COMMON_RUN) {
        /* 3. 共享栈任务先备份自己的栈 */
        if (s_TaskScheduler.pCurrent->pBackupStack != NULL) {
            Task_BackupSharedStack(s_TaskScheduler.pCurrent);
        }
        /* 4. 让出 CPU，寻找下一个可运行任务 */
        Task_JumpNext(&s_TaskScheduler);
    }

    /* 5. 被唤醒后从 setjmp 之后继续执行 */
    TASK_ENABLE_IRQ();
}
\end{lstlisting}

\texttt{Task\_ConditionWait} 的结构完全相同，只是唤醒条件从“时间到”换成“条件变量的标志被置位”。这两段代码共享同一个状态机模板：置 SUSPEND → 保存现场 → 备份栈 → 让出 CPU → 恢复时继续。

由于现场保存在 \texttt{jmp\_buf} 中，任务函数里不需要手动记录：

\begin{lstlisting}[style=c]
void TaskFunc1(int arg)
{
    int cnt = 5;               /* 暂停期间保存在独立栈中 */
    while (cnt--) {
        Task_Wait(1000);       /* 挂起 1 秒，之后从下一行继续 */
        printf("TaskFunc1 loop\n");
    }
}
\end{lstlisting}

这就是“状态机简化”的核心：用户把任务写成顺序结构，编译器生成的调用栈负责记忆执行位置，调度器只负责在合适的时候切换现场。

% ---------------------------------------------------------------------
\section{调度器如何选择下一个任务}

负责选择任务的是 \texttt{Task\_JumpNext}，它的执行过程可以拆成三步：

\begin{enumerate}
\item 扫描唤醒条件：遍历环形任务链表，把定时到点或条件满足的任务从 \texttt{SUSPEND} 置为 \texttt{READY}；
\item 消费条件：清除已被就绪任务“看到”的条件标志；
\item 选择可运行任务：从当前任务的下一个任务开始，找到第一个 \texttt{READY} 任务（或未启动的 \texttt{INIT} 任务），通过 \texttt{longjmp} 切换到它。
\end{enumerate}

\begin{lstlisting}[style=c]
static void Task_JumpNext(TaskScheduler_T *pObj)
{
    while (1) {
        TaskControlBlock_T *p = pObj->pTaskList;

        /* 1. 扫描：把到点/条件满足的任务置为 READY */
        do {
            if (p->state != TASK_STATUS_DELETED &&
                p->state != TASK_STATUS_INIT) {
                if ((p->nextRunTime != 0 &&
                     pObj->pfnGetTime() > p->nextRunTime) ||
                    (p->pWaitCondition != NULL &&
                     p->pWaitCondition->flag == 1)) {
                    p->state = TASK_STATUS_READY;
                }
            }
            p = p->pNext;
        } while (p != pObj->pTaskList);

        /* 2. 找到 READY/INIT 任务并切换 */
        p = pObj->pCurrent->pNext;
        do {
            if (p->state == TASK_STATUS_READY) {
                pObj->pCurrent = p;
                pObj->pCurrent->state = TASK_STATUS_RUNNING;
                longjmp(pObj->pCurrent->env, TASK_COMMON_RUN);
            } else if (p->state == TASK_STATUS_INIT) {
                s_TaskScheduler.pCurrent = p;
                longjmp(s_TaskScheduler.envMain, TASK_MAIN_START_TASK);
            }
            p = p->pNext;
        } while (p != pObj->pCurrent->pNext);
    }
}
\end{lstlisting}

注意：如果暂时没有任务处于 \texttt{READY}，\texttt{Task\_JumpNext} 会继续扫描，直到某个任务到点。因此本章示例是一个“时间驱动”的调度器：当前任务调用 \texttt{Task\_Wait} 后，调度器会一直检查所有挂起任务的唤醒时间。

\texttt{Task\_Start} 使用 \texttt{envMain} 保存“调度器自己”的现场，并用 \texttt{TASK\_MAIN\_*} 系列常量区分返回原因：

\begin{lstlisting}[style=c]
int Task_Start(void)
{
    int ret = setjmp(s_TaskScheduler.envMain);

    if (ret == TASK_MAIN_INIT) {
        /* 首次启动：从第一个任务开始 */
        s_TaskScheduler.pCurrent = s_TaskScheduler.pTaskList;
        longjmp(s_TaskScheduler.envMain, TASK_MAIN_START_TASK);
    } else if (ret == TASK_MAIN_START_TASK) {
        if (s_TaskScheduler.pCurrent->stackType == TASK_STACK_UNIQUE) {
            /* 独立栈：直接切换栈指针后运行入口函数 */
            TASK_SET_STACK(s_TaskScheduler.pCurrent->stackTop);
            s_TaskScheduler.pfnRunTask(TCB_MAGIC_WORD);
        } else {
            Task_StartShared();   /* 共享栈：需要保护跳转栈空间 */
        }
    }
    return 0;
}
\end{lstlisting}

其中 \texttt{TASK\_SET\_STACK} 是一个汇编级的栈指针切换宏。当前代码示例提供的是 x86\_64 版本：

\begin{lstlisting}[style=c]
#define TASK_SET_STACK(p) \
    __asm__ volatile("mov %0, %%rsp" : : "r" (p) : "memory");
\end{lstlisting}

移植到 Cortex-M 时，应改成修改 MSP 与 PSP 寄存器（如 \texttt{\_\_set\_MSP}），或者使用每个架构提供的上下文切换汇编。

% ---------------------------------------------------------------------
\section{独立栈与共享栈的调度差异}

\subsection{独立栈任务}

独立栈任务的切换不涉及栈内容拷贝。启动流程是：

\begin{enumerate}
\item \texttt{Task\_Start} 用 \texttt{TASK\_SET\_STACK} 把栈指针切到任务栈顶；
\item 任务正常执行；
\item 任务调用 \texttt{Task\_Wait} 时，\texttt{setjmp} 把现场保存在自己的 \texttt{env} 中，此时栈指针指向它自己的栈；
\item 任务被再次选中时，\texttt{longjmp} 恢复 \texttt{env}，栈指针回到自己的栈顶，局部变量仍然完好。
\end{enumerate}

\subsection{共享栈任务}

共享栈任务运行在调度器的主栈上，因此切换时必须先“把主栈备份出去”，等它再次被选中时再“搬回来”。代码中对应 \texttt{Task\_StartShared}、\texttt{Task\_BackupSharedStack} 和 \texttt{Task\_RestoreSharedStack}：

\begin{lstlisting}[style=c]
static void Task_StartShared(void)
{
    /* 预留一块安全区，防止恢复共享栈时把跳转现场压坏 */
    volatile char reserve[TASK_SHARED_JMP_STACK_SIZE];
    reserve[0] = 0;

    s_TaskScheduler.pCurrent->stackTop =
        s_TaskScheduler.pfnGetStackTop();

    if (setjmp(s_TaskScheduler.pCurrent->env) == TASK_COMMON_INIT) {
        s_TaskScheduler.pfnRunTask(TCB_MAGIC_WORD);
    }
}
\end{lstlisting}

\texttt{reserve} 数组的存在很关键：它保证了从“备份共享栈”到“longjmp 回调度器主循环”这段过程中，跳转本身需要的栈帧不会被任务数据覆盖。

% ---------------------------------------------------------------------
\section{栈溢出保护}

独立栈虽然简单，但栈大小完全由用户指定，一旦给少了就会越界写入相邻变量。本章代码在任务入口处使用了一个“魔法字”做粗略检测：

\begin{lstlisting}[style=c]
static void Task_RunEntry(uint32_t magic)
{
    TASK_ENABLE_IRQ();
    s_TaskScheduler.pCurrent->state = TASK_STATUS_RUNNING;
    s_TaskScheduler.pCurrent->pfnEntry(s_TaskScheduler.pCurrent->entryArg);

    /* 如果栈被破坏，magic 参数不可能还等于 TCB_MAGIC_WORD */
    while (magic != TCB_MAGIC_WORD);

    s_TaskScheduler.pCurrent->state = TASK_STATUS_DELETED;
}
\end{lstlisting}

\texttt{magic} 是任务启动时从调用方一路传进来的参数，正常返回时它的值必须仍然是 \texttt{0xDEADBEEF}。如果任务执行期间栈溢出把参数覆盖了，代码会停在 \texttt{while} 循环里，让开发者立刻意识到栈太小。这只是防御性检测，正确做法仍然是在开发阶段给每个任务留足余量，并通过“压栈填充 + 水位检查”工具测量实际最大栈用量。

% ---------------------------------------------------------------------
\section{两个任务的配示例合}

\texttt{main.c} 演示了“生产-通知-等待”模式：\texttt{TaskFunc1} 每 1 秒通知一次条件变量， \texttt{TaskFunc2} 等待条件变量后打印消息：

\begin{lstlisting}[style=c]
TaskCondition_T g_Condition;

void TaskFunc1(int arg)
{
    int cnt = 5;
    while (cnt--) {
        Task_Wait(1000);                    /* 挂起 1 秒 */
        printf("TaskFunc1 loop\n");
        Task_ConditionNotify(&g_Condition); /* 唤醒等待者 */
    }
}

void TaskFunc2(int arg)
{
    while (1) {
        Task_ConditionWait(&g_Condition);   /* 没有通知就一直挂起 */
        printf("TaskFunc2 got condition\n");
    }
}

int main(void)
{
    Task_Init(GetTimerMs);
    Task_ConditionInit(&g_Condition);

    Task_CreateWithStack(TaskFunc1, "task1", s_Stack1, sizeof(s_Stack1), 10);
    Task_CreateWithStack(TaskFunc2, "task2", s_Stack2, sizeof(s_Stack2), 20);
    Task_Create(TaskFunc1, "task3", 30);    /* 共享栈任务 */

    Task_Start();                           /* 不会返回 */
    return 0;
}
\end{lstlisting}

当 \texttt{TaskFunc2} 调用 \texttt{Task\_ConditionWait} 时，它会被挂起并让出 CPU；1 秒后 \texttt{TaskFunc1} 调用任务通知函数 \texttt{Task\_ConditionNotify} 把条件置 1，下一次 \texttt{Task\_JumpNext} 扫描时就把 \texttt{TaskFunc2} 置为 READY 并恢复它。

% ---------------------------------------------------------------------
\section{与其他框架的架构差异}

理解一个调度器的设计，最好的方式不是孤立地看它的代码，而是把它放在“备选方案”的坐标系中比较。本章的协程调度器位于裸机状态机与完整 RTOS 之间的“中间地带”——它比第~\ref{ch:simple-task}~章的 \texttt{simple\_task} 多了一整套栈管理机制，能支持任务在任意深度暂停后恢复；但它又远不如 FreeRTOS 那样复杂，没有抢占、没有优先级、没有信号量和队列，只是一个“能让函数挂起和恢复”的轻量框架。

\subsection{与 simple\_task 状态机的对比}

把本章调度器与第~\ref{ch:simple-task}~章的 \texttt{simple\_task} 放在一起，可以更清楚地看到两种设计的分界线，如表~\ref{tab:coro-vs-simple}~所示。

\begin{table}[htbp]
	\centering
	\small
	\caption{增强版协程调度器与 simple\_task 的对比}
	\label{tab:coro-vs-simple}
	\begin{tabularx}{\textwidth}{l X X}   % 后两列自动换行
		\toprule
		对比项 & simple\_task & 本章协程调度器 \\
		\midrule
		状态集合 & WAIT、READY、RUN、DISABLE、ONCE & INIT、READY、RUNNING、SUSPEND、DELETED \\
		任务能否中途挂起 & 不能，回调每次运行后返回 & 能，可在任意函数内暂停后恢复 \\
		暂停方式 & 返回延时值，下次从回调入口重跑 & setjmp 保存现场，longjmp 恢复现场 \\
		栈管理 & 任务共用主循环栈 & 独立栈，或共享栈+备份拷贝 \\
		状态迁移位置 & Task\_Tick 与 Task\_RunTask 显式迁移 & Task\_JumpNext + longjmp 集中迁移 \\
		定时来源 & 1ms 定时器中断计数 & 外部毫秒时间回调 \\
		就绪队列 & FIFO 队列（可关闭） & 环形任务链表轮询扫描 \\
		优点 & 纯 C、RAM 小、可移植性高 & 代码写法接近顺序执行，适合复杂流程 \\
		代价 & 复杂流程必须拆成状态机 & 需要栈空间和汇编支持，移植成本高 \\
		\bottomrule
	\end{tabularx}
\end{table}

从状态机角度看，两者都要回答同一个问题：一个任务什么时候能跑、跑完把状态切成什么。区别在于：

\begin{enumerate}
	\item simple\_task 是“数据驱动的显式状态机”：任务状态存在 \texttt{Task\_T.systemState} 中，定时器中断里的点滴函数 \texttt{Task\_Tick} 负责 WAIT→READY，主循环里的 \texttt{Task\_RunTask} 负责 READY→RUN→WAIT/DISABLE。用户看到的状态机是“可见”的；
	\item 本章协程是“现场驱动的隐式状态机”：任务暂停点就是状态机的一个分支，状态保存在 \texttt{jmp\_buf} 和任务栈中。用户看到的是顺序代码，切换细节被封装进 \texttt{Task\_Wait} 与 \texttt{Task\_ConditionWait}。
\end{enumerate}

在 \texttt{simple\_task} 中，一个按键长按+双击+重复的复杂流程，通常要拆成回调里的一堆用户状态；在本章协程中，可以直接写成：

\begin{lstlisting}[style=c]
	void KeyFlowTask(int arg)
	{
		while (1) {
			Task_Wait(Key_GetPressMs());     /* 等待按键事件 */
			if (Key_IsDouble()) {
				HandleDoubleClick();
				Task_Wait(Key_ReleaseDone());
			} else if (Key_IsLong()) {
				HandleLongPress();
			}
		}
	}
\end{lstlisting}

这不是说协程一定优于 \texttt{simple\_task}。对于大量短小的周期任务（每 1ms 或 10ms 翻转一个 LED、轮询一个按键），\texttt{simple\_task} 更轻、更省内存、更好调试；对于包含“等待外部事件→处理→继续等待”的复杂协议流程，有栈协程能让代码结构更接近人的思维。

\subsection{与RTOS框架的对比}

如果把本章的协程调度器与一款完整的嵌入式 RTOS（如 FreeRTOS、RT-Thread Nano 或 Zephyr）放在一起比较，可以更清楚地看到“有栈协程”在资源受限场景下的定位，如表~\ref{tab:coro-vs-rtos}~所示。

\begin{table}[htbp]
	\centering
	\small
	\caption{增强版协程调度器与 RTOS 的对比}
	\label{tab:coro-vs-rtos}
	\begin{tabularx}{\textwidth}{l X X}
		\toprule
		对比项 & 本章协程调度器 & 完整 RTOS（如 FreeRTOS） \\
		\midrule
		调度方式 & 协作式（任务主动让出 CPU） & 抢占式（支持时间片轮转与优先级抢占） \\
		任务切换 & setjmp/longjmp，仅保存必要寄存器 &  PendSV 中断 + 硬件上下文切换，保存全部寄存器 \\
		栈管理 & 独立栈或共享栈+备份，用户手动分配 & 每个任务独立栈，通常由创建函数动态分配 \\
		中断响应 & 任务切换时关中断，中断延迟受临界区影响 & 中断嵌套，PendSV 将上下文切换延迟到最低优先级 \\
		任务间通信 & Task\_ConditionWait/Notify（条件变量） & 队列、信号量、互斥量、事件组等丰富机制 \\
		内存占用 & 数 KB（TCB + 任务栈） & 数十 KB（内核对象 + 多任务栈） \\
		移植成本 & 需实现 TASK\_SET\_STACK 与时间源 & 需实现时钟节拍中断与上下文切换汇编 \\
		调试难度 & 中（setjmp 跳转使调用栈不连续） & 中高（抢占导致执行顺序不确定，需借助 Trace 工具） \\
		典型应用 & 复杂协议流程、事件驱动型任务 & 多优先级实时系统、复杂外设驱动管理 \\
		\bottomrule
	\end{tabularx}
\end{table}

从表中可以看到，RTOS 的核心优势在于：

\begin{enumerate}
	\item 抢占式调度：高优先级任务不会被低优先级任务阻塞（除非使用互斥量导致的优先级反转）。对于需要严格实时响应的场景（如电机控制回路、通信协议栈的超时处理），抢占调度是不可或缺的；
	\item 丰富的同步原语：信号量、队列、互斥量等机制为标准组件，多任务协作时不必自己实现条件变量或消息缓冲；
	\item 成熟的生态：FreeRTOS 等有大量现成移植、驱动适配和调试工具（如 Tracealyzer、SystemView），能够降低大型项目的长期维护成本。
\end{enumerate}

\begin{anybox}{工程观点：特斯勒定理与任务边界}
特斯勒定理（复杂度守恒定律）指出：一个系统中总有一部分固有复杂度无法消除，它只能在框架、应用代码与使用者之间转移，最终必须由某一方承担。因此，RTOS 并不一定是一个很坏的框架；当系统需要网络协议栈或文件系统协议栈时，使用成熟 RTOS 往往比自行补齐调度、同步和协议适配机制更稳妥。

如果担心抢占式 RTOS 带来的执行时序不确定性，可以只设置两到三个职责明确、粒度较粗的任务，用简单的信号量完成任务间的同步或通知，而不必把每个传感器、UI 组件等都拆成独立任务。对于这些粗粒度任务内部的顺序流程，可以采用本章基于 \texttt{setjmp}/\texttt{longjmp} 的有栈协程，也可以采用第~\ref{ch:simple-task}~章介绍的 \texttt{simple\_task} 协程式状态机。需要注意的是，如果既把 RTOS 任务拆得过细，又在每个任务中叠加协程或状态机，复杂度并不会消失，反而会转移到任务边界、同步关系、栈空间和调试过程之中。因此，选择框架的关键不是排斥 RTOS，而是控制任务数量，并让每一种调度模型只承担它最适合表达的那部分复杂度。
\end{anybox}

而本章协程调度器相对于 RTOS 的“轻量”并非简单的功能裁剪，而是换了一种思路：

\begin{enumerate}
	\item 协作式调度的确定性：任务只在显式调用 \texttt{Task\_Wait} 或返回时才切换，因此共享数据的临界区天然短小。而 RTOS 的抢占切换可能发生在任意指令边界，必须用信号量或关中断保护共享资源；
	\item 无需修改硬件寄存器：在 Cortex-M 等平台上，RTOS 必须实现 PendSV 中断和 \texttt{svc} 指令调用，涉及对 \texttt{CONTROL} 寄存器、PSP 指针的操作；而本章调度器只需提供 \texttt{TASK\_SET\_STACK} 宏，不涉及中断向量表或硬件特权级切换，移植量明显更小；
	\item 内存占用可精确控制：RTOS 的每个任务栈通常由 \texttt{pvPortMalloc} 从堆中分配，且栈大小需按最坏情况预分配；本章的独立栈使用静态数组，栈溢出检测也更容易和编译器配合；
	\item 无内核对象管理开销：RTOS 的信号量、队列等内核对象需要额外的 RAM 和初始化代码，任务数量增多时创建和删除的开销逐渐不可忽略；本章调度器的任务链表和条件变量均为静态结构，没有动态内存分配。
\end{enumerate}

因此，是否使用 RTOS 取决于项目的具体需求：

\begin{itemize}
	\item 适合使用 RTOS 的场景：系统包含多个优先级明确的任务（如 1kHz 控制循环、UI 刷新、通信协议栈共存），或需要标准化的任务间通信机制（队列、信号量），且 MCU 的 RAM/Flash 足够（通常要求 32KB+ RAM、128KB+ Flash）；
	\item 适合使用本章协程调度器的场景：系统主要为“等待事件→处理→继续等待”的流程型任务（如命令解析、传感器采集序列、人机交互流程），任务间仅有简单的通知关系，且 MCU 资源较紧张（RAM < 16KB）；
	\item 兼顾两者的混合方案：也可以在 RTOS 的一个任务里运行本章协程调度器——把协程作为 RTOS 线程内部的“子任务”，既利用 RTOS 的抢占和通信能力，又在单个线程内部用协程简化多个顺序子流程的编写。
\end{itemize}

对于单片机开发而言，选择调度器不是技术崇拜，而是资源与需求的精确匹配。本章的目的不是造一个“迷你 RTOS”，而是在“裸机状态机”和“完整 RTOS”之间提供一个中间选项——它比状态机更适合表达复杂流程，又比 RTOS 更轻量、更可控。读者应结合自己项目的 RAM 预算、任务结构和实时性要求，在三种方案之间做出理性的选择。

% ---------------------------------------------------------------------
\section{移植要点与限制}

\subsection{移植时至少需要提供}

\begin{enumerate}
\item \texttt{TaskGetTime\_f}：可用的毫秒时间源（如 SysTick、HAL\_GetTick）；
\item \texttt{TASK\_DISABLE\_IRQ} 与 \texttt{TASK\_ENABLE\_IRQ}：调度临界区的中断开关；
\item \texttt{TASK\_SET\_STACK}：修改当前栈指针的底层实现，需按架构编写；
\item 栈区必须满足 16 字节对齐，并保证不与全局变量或堆重叠。
\end{enumerate}

\subsection{使用限制}

\begin{itemize}
\item 不要在任务入口外的普通函数里调用 \texttt{Task\_Wait} 后，还依赖该函数的 \texttt{jmp\_buf} 在其返回后继续生效；
\item 共享栈任务禁止在深层嵌套的递归函数中挂起，否则备份区会溢出；
\item \texttt{longjmp} 不能跨越已释放的栈帧，因此要等待任务结束后再回收它的 TCB；
\item 当前版本是协作式调度：任务自己不 \texttt{Task\_Wait}，其他任务就得不到 CPU，不适合硬实时或抢占需求。
\end{itemize}

% ---------------------------------------------------------------------
\section{实验报告：协程调度器移植与验证}

\subsection{实验基本信息}

MCU 型号：\LabFillLine[3.5cm] \quad 工作主频：\LabFillLine[3.5cm]

毫秒计时源：\LabFillLine[3.5cm] \quad 栈对齐方式：\LabFillLine[3.5cm]

独立栈任务数：\LabFillLine[3.5cm] \quad 共享栈任务数：\LabFillLine[3.5cm]

\subsection{实验目的}
\begin{enumerate}
\item 在 Cortex-M 上实现毫秒时间源、中断临界区宏和栈指针切换宏；
\item 验证独立栈与共享栈任务的创建、挂起（Task\_Wait）和恢复过程；
\item 观察任务运行日志，测量并分析两种栈模式的 RAM 占用与切换开销。
\end{enumerate}

\subsection{实验原理}
调度器使用 \texttt{setjmp}/\texttt{longjmp} 保存和恢复任务执行现场（寄存器、栈指针、程序计数器）。独立栈任务切换时仅需修改栈指针；共享栈任务在挂起时将当前运行栈的内容备份到私有缓冲区，恢复时再拷贝回去。调度器通过环形任务链表轮询唤醒条件，选择就绪任务并跳转执行。

\subsection{硬件连接}

\LabPasteBox{协程调度器硬件连接}{粘贴开发板、调试串口（或 LED）的连接图片，标注电源、调试接口和指示灯。}{fig:ch9-coro-hardware}

\subsection{关键代码}

\LabPasteBox{时间源、栈切换宏与任务创建代码}{粘贴 \texttt{GetTimerMs}、\texttt{TASK\_SET\_STACK} 宏实现，以及独立栈与共享栈任务的创建代码，并标注栈大小和对齐检查。}{fig:ch9-coro-code}

\subsection{实验步骤}
\begin{enumerate}
\item 提供毫秒级计时函数（如 SysTick 中断计数），并实现 \texttt{TASK\_DISABLE\_IRQ} 与 \texttt{TASK\_ENABLE\_IRQ} 临界区宏。
\item 根据目标架构实现 \texttt{TASK\_SET\_STACK} 宏（Cortex-M 可修改 PSP，x86 可修改 RSP），并确保栈指针 16 字节对齐。
\item 配置 \texttt{task\_scheduler\_config.h}，设定任务数量、共享栈备份区大小等参数。
\item 创建两个独立栈任务：任务 A 每 1 秒翻转 LED，任务 B 每 2 秒通过串口打印日志。
\item 创建一个共享栈任务，循环调用 \texttt{Task\_Wait} 多次，观察局部变量是否在恢复后保持不变。
\item 通过串口输出任务运行时间戳和状态切换日志，记录各任务实际最大栈用量（可通过栈水印检测）。
\item 依次注释掉独立栈和共享栈任务，对比两种模式下的 RAM 占用（查看 MAP 文件）。
\end{enumerate}

\subsection{实验现象与数据记录}

\begin{table}[htbp]
\centering
\caption{协程任务运行记录}
\label{tab:ch9-coro-record}
\begin{tabularx}{\textwidth}{p{2.6cm}p{2.8cm}p{2.6cm}X}
\toprule
任务名 & 栈模式 & 周期（ms） & 运行日志/状态 \\
\midrule
\LabFillLine[2.2cm] & \LabFillLine[2.2cm] & \LabFillLine[2.2cm] & \LabFillLine[2.2cm] \\
\LabFillLine[2.2cm] & \LabFillLine[2.2cm] & \LabFillLine[2.2cm] & \LabFillLine[2.2cm] \\
\LabFillLine[2.2cm] & \LabFillLine[2.2cm] & \LabFillLine[2.2cm] & \LabFillLine[2.2cm] \\
\bottomrule
\end{tabularx}
\end{table}

\LabPasteBox{协程调度器运行结果}{粘贴串口打印的任务切换日志、LED 翻转波形截图或逻辑分析仪捕获的时序图。}{fig:ch9-coro-result}

\subsection{结果分析}

独立栈任务实际最大栈用量（字节）：\LabFillLine[6cm]

共享栈任务备份区最大使用量（字节）：\LabFillLine[6cm]

切换耗时对比（独立栈 vs. 共享栈）：\LabFillLine[6cm]

移植过程中遇到的对齐或栈溢出异常：\LabFillLine[6cm]

\subsection{实验总结}

\noindent\LabFillLine[\linewidth]\par
\noindent\LabFillLine[\linewidth]\par
\noindent\LabFillLine[\linewidth]

% ---------------------------------------------------------------------
\section{小结}

本章涉及的栈切换、现场保存、中断与任务上下文，都依赖于对 CPU 内核的理解。如果希望深入掌握 Cortex-M 体系结构与嵌入式 C 的底层细节，以下书目值得一读：

\begin{itemize}
	\item 《ARM Cortex-M3 与 Cortex-M4 权威指南》（第 3 版） —— Joseph Yiu 著，清华大学出版社。Cortex-M 内核的“圣经”，从处理器模型、异常处理、内存映射到浮点单元，覆盖了协程调度器移植所需的所有硬件知识（MSP/PSP、PendSV、堆栈对齐等）。
	\item 《C 嵌入式编程设计模式》（英文名 \textit{Design Patterns for Embedded Systems in C}） —— Bruce Powel Douglass 著，机械工业出版社。以状态机、任务调度、并发模式等实用视角讲解嵌入式 C 的设计取舍，可与本章的调度器实现互相印证。
	\item 《嵌入式系统设计：CPS 与物联网物理世界》 —— 陈海波、夏虞斌等编著，高等教育出版社。从系统软件层面讲解中断、异常、实时调度和虚拟化，适合希望从“裸机”迈向“系统软件”视角的读者。
\end{itemize}

这些书不需要通读，建议在遇到本章移植疑惑（如“为什么TASK\_SET\_STACK要改PSP”、“setjmp在ARM上需要保存哪些寄存器”）时，按需查阅对应章节。

本章用 \texttt{setjmp} 与 \texttt{longjmp} 实现了一个有栈协程调度器，核心收获有三点：

\begin{enumerate}
	\item setjmp 与 longjmp 是“现场保存和恢复”原语：它把一个普通 C 函数变成可以挂起、恢复的协程；
	\item 栈是任务暂停时的“记忆”：自动变量和调用链都存放在栈上，所以任务需要独立栈，或用备份区把共享栈内容保存下来；
	\item 状态机可以被“藏”进上下文切换：\texttt{simple\_task} 用枚举+循环显式表达状态，本章调度器用 \texttt{jmp\_buf} 隐式记住执行位置。两者不是替代关系，而是 “轻量周期任务”与“复杂阻塞流程”两种场景下的不同工具。
\end{enumerate}

如果读者需要在资源更受限的 51 或 STM8 上实现类似效果，可以只保留“独立栈+状态机”的子集，把汇编级栈切换换成纯状态机；如果追求代码最接近顺序思维，则可以在此基础上继续引入 Protothreads、协程宏或完整 RTOS 的线程模型。

% ---------------------------------------------------------------------
\begin{problemset}
\item \texttt{setjmp} 保存了什么现场？\texttt{longjmp} 恢复现场后，哪些局部变量不能假设其值仍然有效？
\item 为什么每个协程任务需要“独立栈”？如果改用共享栈，切换时为什么必须备份栈内容？
\item 请画出本章协程调度器从 READY→RUNNING→SUSPEND→READY 的状态迁移，并说明哪些操作发生在中断里。
\item 与 \texttt{simple\_task} 的显式状态机相比，协程把状态“藏”在哪里？这带来什么优点和调试难点？
\item 协程任务栈发生溢出时会破坏什么？你能设计一个“栈水印”来检测最大栈用量吗？
\item Cortex-M 的 MSP 与 PSP 有什么区别？切换任务栈时需要修改哪一个寄存器？
\item 查阅资料，使用 C语言和 ARM 汇编实现 \texttt{setjmp} 和 \texttt{longjmp} 函数。
\begin{enumerate}
	\item 在 ARM 架构中，\texttt{jmp\_buf} 通常需要保存哪些寄存器？为什么？
	\item 查阅 \texttt{newlib}、\texttt{klibc} 或 \texttt{Barebox} 等开源项目中 \texttt{arch/arm/setjmp.S} 的实现[reference:0]，说明 \texttt{setjmp} 的核心汇编指令（如 \texttt{stmia}）是如何保存上下文的。
	\item \texttt{longjmp} 如何恢复上下文并跳转到 \texttt{setjmp} 的返回位置？为什么 \texttt{longjmp} 需要确保返回值不为 0？
	\item ARM 与 Thumb 指令集下 \texttt{setjmp} 的实现有何不同？请说明 Thumb 模式下如何处理高寄存器（\texttt{r8}-\texttt{r15}）的保存与恢复。
	\item 若在 Cortex-M 上实现协程调度器，除了 \texttt{setjmp}/\texttt{longjmp} 外，还需要额外处理什么（如 MSP 与 PSP 的切换）？
\end{enumerate}
\end{problemset}
